\section{From finite tensor powers to absorption}\label{sec:application}

\begin{proof}[Proof of Theorem~\ref{thm:finite}]
Proposition~\ref{prop:full-cone} supplies a finite tensor power
$P=D^{\otimes m_0}$ and a unital map $\eta:C\to P$ for which
$\eta(\iota e_{11})$ is full. Proposition~\ref{prop:cone-drop} supplies a
unital map from $I(2,3)$ into a finite tensor power $P^{\otimes s}$.
Associativity of the maximal tensor product identifies the latter with
$D^{\otimes m_0s}$, as required.
\end{proof}

For the application, we first record explicitly why the possibly
nonseparable central-sequence algebra contains a suitable separable
subalgebra. This is the separable-reduction argument of
\cite[Lemma~3.5]{KR}; we include its finite commutator witness.

\begin{lemma}\label{lem:separable}
Every nonzero unital $C^*$-algebra $B$ without characters contains a
separable $C^*$-subalgebra $D$ without characters and with $1_D=1_B$.
\end{lemma}
\begin{proof}
The closed ideal generated by commutators of selfadjoint elements is $B$;
a proper such ideal would give a nonzero unital commutative quotient and
hence a character. Consequently, there are finitely many
$h_i,k_i\in B$ and $b_i,c_i\in B_{\sa}$ such that
\[
 \left\|1_B-\sum_{i=1}^r h_i[b_i,c_i]k_i\right\|<1.
\]
Let $D$ be generated by $1_B$ and all the displayed elements. It is unital
and separable. A character of $D$ would send the sum to zero and the unit
to one, contradicting its contractivity and the strict norm bound.
\end{proof}

\begin{proof}[Proof of Theorem~\ref{thm:main}]
Assume that $F_\omega(A)$ has no characters. This algebra is nonzero and
unital because it contains the diagonal unit of $A$. By
Lemma~\ref{lem:separable}, choose a separable unital characterless
subalgebra $D\subset F_\omega(A)$ with the same unit. By
Theorem~\ref{thm:finite} and Proposition~\ref{prop:kr}(1), there is a
composition of unital $*$-homomorphisms
\[
 I(2,3)\longrightarrow D^{\otimes_{\max}m}
 \longrightarrow F_\omega(A).
\]
Denote its image by $E$. The image is separable and is a quotient of
$I(2,3)$. Every irreducible representation of $I(2,3)$ factors through
evaluation at a point of $[0,1]$: the central algebra $C([0,1])1$ acts
by a character, and sections vanishing at its evaluation point are
annihilated, by approximation using scalar functions vanishing there.
The resulting irreducible dimensions are $2$ at $0$, $3$ at $1$, and
$6$ in the interior. Thus $I(2,3)$ is subhomogeneous and has no
characters. The same properties hold for $E$, since any irreducible
representation or character of $E$ pulls back along the quotient map.
Also $1_E=1_{F_\omega(A)}$. Proposition~\ref{prop:kr}(2) gives
$A\cong A\otmin\Z$.

Conversely, if $A\cong A\otmin\Z$, Proposition~\ref{prop:kr}(2)
supplies a unital subalgebra $E\subset F_\omega(A)$ without characters.
Any character of $F_\omega(A)$ would restrict to a character of $E$,
a contradiction.
\end{proof}

\begin{proof}[Proof of Corollary~\ref{cor:infinite-tensor}]
Because $D$ is nonzero and unital, each spatial connecting map
$x\mapsto x\otimes1_D$ is a unital embedding. Thus $B$ is nonzero, unital,
and separable, and the canonical maps
$\iota_m:D^{\otimes_{\min}m}\hookrightarrow B$ are unital.
Choose the unital homomorphism $\phi$ from Theorem~\ref{thm:finite} and
compose it with the canonical maximal-to-minimal quotient and the
finite-stage inclusion:
\[
 I(2,3)\xrightarrow{\ \phi\ }D^{\otimes_{\max}m}
 \xrightarrow{\ q_m\ }D^{\otimes_{\min}m}
 \xrightarrow{\ \iota_m\ }B.
\]
The quotient $q_m$ is unital, with no nuclearity assumption on $D$.
Let $\theta$ denote this composite and put $E=\theta(I(2,3))$.
Since $\theta$ is unital, $1_E=1_B\ne0$.
The algebra $I(2,3)$ is separable and, as checked above, subhomogeneous
and characterless. Its quotient $E\cong I(2,3)/\ker\theta$ is therefore
separable and subhomogeneous: every irreducible representation of $E$
pulls back to one of $I(2,3)$. A character of $E$ would compose with the
unital quotient map to a character of $I(2,3)$, which is impossible.
These conclusions hold even when $\theta$ has a kernel.

Dadarlat--Toms \cite[Theorem~1.1]{DT}, applied to the unital separable
algebra $D$ and this same-unit subalgebra of $B$, gives
$B\cong B\otmin\Z$. Finally, because $B\ne0$, the map
$z\mapsto1_B\otimes z$ embeds $\Z$ unitally into $B\otmin\Z$.
Transporting this map through the inverse of an isomorphism
$B\cong B\otmin\Z$, which is automatically unital, gives the asserted
unital copy of $\Z$ in $B$.
\end{proof}
